\begin{longtable}{@{}l p{12.5cm}@{}}
  \caption{Pedagogical per-instance answer comparison ($\Delta \geq 75\%$) for the 5 most divergent items: same question, each model's answer and reasoning.}\label{tab:cmp_pedagogical_pos75} \\
  \toprule
  \endfirsthead
  \toprule
  \endhead
  \multicolumn{2}{@{}l}{\textbf{chile\_110} \quad $\Delta = +100\%$} \\
  \textbf{Question} & To work on the concept of hydrostatic pressure, a third-year high school teacher plans a class with the objective for students to apply the concept of density when comparing different bodies and substances. At the start of the class, an activity will be conducted to introduce the concept of density. Which of the following tasks should students perform, after the initial activity mentioned above, in order to address the objective of the class? \newline A. Predict what will happen when three immiscible liquids are poured into a graduated cylinder, and explain why the observed phenomenon occurs \newline B. Determine experimentally the density of a liquid based on the pressure it exerts on the bottom of the container holding it \newline C. Explain how the melting of Antarctic icebergs affects sea level, based on the density of ice and liquid water \newline D. Relate qualitatively and quantitatively the density of different bodies with the same mass and different volume, or the same volume and different mass \\
  \textbf{Gold} & D \\
  \midrule
  \textbf{Kimi-K2.6} & answer A \\
  & The user wants me to evaluate a multiple-choice question and answer with ONLY the letter of the correct answer. \newline Let me analyze the question carefully: \newline Context: A third-year high school teacher plans a class on hydrostatic pressure. The objective is for students to apply the concept of density when comparing different bodies and substances. \newline At the start: an activity to introduce the concept of density. \newline Question: Which task should students perform AFTER the initial activity to address the objective of the class? \newline Key points: \newline \textit{[\ldots\ truncated \ldots]} \newline One more check: Is there any reason C would be correct? Explaining sea level rise based on density. This is an application, but it's not directly related to hydrostatic pressure classwork, and it's more about buoyancy. Also, it doesn't involve comparing multiple bodies/substances in an experimental/classroom context. \newline So final answer: A. \\
  \textbf{Innospark-1T} & answer D \\
  & Let me analyze this question step by step. \newline The question asks: To work on the concept of hydrostatic pressure, a third-year high school teacher plans a class with the objective for students to apply the concept of density when comparing different bodies and substances. At the start of the class, an activity will be conducted to introduce the concept of density. Which of the following tasks should students perform, after th \newline Let me consider each option: \newline - Option A: Predict what will happen when three immiscible liquids are poured into a graduated cylinder, and explain why the observed phenomenon occurs \newline - Option B: Determine experimentally the density of a liquid based on the pressure it exerts on the bottom of the container holding it \newline - Option C: Explain how the melting of Antarctic icebergs affects sea level, based on the density of ice and liquid water \newline \textit{[\ldots\ truncated \ldots]} \newline This option most directly addresses the stated objective by having students apply the concept of density to compare different bodies and substances, establishing both the conceptual understanding and mathematical relationship needed for subsequent work with hydrostatic pressure. \newline Based on this analysis, the answer is D. \\
  \specialrule{1.2pt}{8pt}{8pt}
  \multicolumn{2}{@{}l}{\textbf{chile\_666} \quad $\Delta = +100\%$} \\
  \textbf{Question} & A teacher will begin working on the unit 'Economic models and their impact on the territory', with the goal of having her high school students analyze the transformations experienced by the geographical space as a result of the relationships between economic, demographic, and spatial processes. Therefore, she has decided to start the class by recalling the meaning of the term geographical space, through a group brainstorming session. In the responses given by the students, those that characterized the geographical space exclusively from its physical or natural component prevailed, attributing to it an attribute of stability in the medium and long term. Considering the didactic principles that guide the teaching of geography, which one of them is essential to work on in order to expand the concept of space that this group of students has? \newline A. Connection \newline B. Location \newline C. Interdependence \newline D. Causality \\
  \textbf{Gold} & C \\
  \midrule
  \textbf{Kimi-K2.6} & answer A \\
  & The user wants me to answer a multiple-choice question about didactic principles in geography teaching. \newline Let me analyze the question: \newline - Unit: 'Economic models and their impact on the territory' \newline - Goal: analyze transformations experienced by geographical space as a result of relationships between economic, demographic, and spatial processes \newline - Current student understanding: geographical space characterized exclusively from its physical/natural component, attributing stability in medium/long term \newline - Need: expand the concept of space beyond just physical/natural stability \newline \textit{[\ldots\ truncated \ldots]} \newline And the principle of connection is indeed the one that integrates the components of geographical space. \newline Final answer: A. \\
  \textbf{Innospark-1T} & answer C \\
  & Let me analyze this question step by step. \newline The question asks: A teacher will begin working on the unit 'Economic models and their impact on the territory', with the goal of having her high school students analyze the transformations experienced by the geographical space as a result of the relationships between economic, demographic, and spatial processes. Therefore, she has decided to start the class by recalling the \newline Let me consider each option: \newline - Option A: Connection \newline - Option B: Location \newline - Option C: Interdependence \newline \textit{[\ldots\ truncated \ldots]} \newline Final Answer: C. \newline Based on this analysis, the answer is C. \\
  \specialrule{1.2pt}{8pt}{8pt}
  \multicolumn{2}{@{}l}{\textbf{chile\_574} \quad $\Delta = +100\%$} \\
  \textbf{Question} & A fourth-grade teacher is planning a class in which she will work on comparing decimal numbers. To start, she will ask the students to compare the following numbers: 1.37 with 1.4 and 1.25 with 1.5; for this purpose, she is looking for a type of representation that will help them perform the exercise and understand the content. Which of the following activities would meet the requirement that the teacher has set? \newline A. Write the quantities as hundredths \newline B. Express the quantities in fractional form \newline C. Represent the quantities with concrete materials \newline D. Illustrate the quantities in everyday situations \\
  \textbf{Gold} & A \\
  \midrule
  \textbf{Kimi-K2.6} & answer C \\
  & The user wants me to answer a multiple-choice question about teaching decimal comparison to fourth-grade students. I need to choose the best activity that helps students compare decimals like 1.37 vs 1.4 and 1.25 vs 1.5. \newline Let me analyze each option: \newline A. Write the quantities as hundredths - This would mean writing 1.4 as 1.40 and 1.5 as 1.50. This is a good mathematical strategy for comparing decimals, but the question asks for a "type of representation" that helps them perform the exercise and understand the content. Writing as hundredths is more of a procedural step. \newline B. Express the quantities in fractional form - Converting to fractions (like 137/100 vs 14/10 or 140/100) might actually make it harder for fourth graders to compare, and it's not the most intuitive representation for understanding decimal magnitude. \newline C. Represent the quantities with concrete materials - Using base-ten blocks, decimal squares, number lines, or other manipulatives would provide a concrete representation. This is often recommended in math education for conceptual understanding. For example, showing that 1.4 is 1 whole and 4 tenths, while 1.37 is 1 whole and 3 tenths and 7 hundredths - visually seeing that 4 tenths \textgreater{} 3 tenths and 7 hundredths. \newline D. Illustrate the quantities in everyday situations - Real-world contexts (like money) could help, but money might be confusing since \$1.4 would be \$1.40, and students might not naturally think of \$1.4 as valid notation. Everyday situations are good for application but may not be the best "representation" for the initial comparison task. \newline \textit{[\ldots\ truncated \ldots]} \newline Therefore, C seems pedagogically superior for understanding. \newline I'm going to go with C. \\
  \textbf{Innospark-1T} & answer A \\
  & Let me analyze this question step by step. \newline The question asks: A fourth-grade teacher is planning a class in which she will work on comparing decimal numbers. To start, she will ask the students to compare the following numbers: 1.37 with 1.4 and 1.25 with 1.5; for this purpose, she is looking for a type of representation that will help them perform the exercise and understand the cont \newline Let me consider each option: \newline - Option A: Write the quantities as hundredths \newline - option B: Express the quantities in fractional form \newline - option C: Represent the quantities with concrete materials \newline \textit{[\ldots\ truncated \ldots]} \newline This is the most direct and effective representation because it makes the place value structure explicit and eliminates the visual confusion that leads to common misconceptions about decimal comparison. \newline Based on this analysis, the answer is A. \\
  \specialrule{1.2pt}{8pt}{8pt}
  \multicolumn{2}{@{}l}{\textbf{chile\_714} \quad $\Delta = +100\%$} \\
  \textbf{Question} & A teacher will start working with her 6th grade students based on the following plan: 6th Grade Unit: The political organization of Chile and the basic principles established by the Constitution. Learning Objective: Explain some fundamental elements of the democratic organization of Chile, including the division of state powers and their representation through elected officials (councilors, mayors, deputies, senators, and President). Evaluation indicator: Recognize the main characteristics that a representative democratic regime must have. Which of the following contents is essential for the students to manage in order to achieve the expected performance? \newline A. Citizen participation \newline B. Rights and duties \newline C. Type of State \newline D. Method of appointing authorities \\
  \textbf{Gold} & A \\
  \midrule
  \textbf{Kimi-K2.6} & answer D \\
  & The user wants me to answer a multiple-choice question about a 6th grade unit on the political organization of Chile and democratic principles. \newline Let me analyze the question: \newline Context: \newline - Unit: The political organization of Chile and the basic principles established by the Constitution. \newline - Learning Objective: Explain some fundamental elements of the democratic organization of Chile, including the division of state powers and their representation through elected officials (councilors, mayors, deputies, senators, and President). \newline - Evaluation indicator: Recognize the main characteristics that a representative democratic regime must have. \newline \textit{[\ldots\ truncated \ldots]} \newline Also, in many standardized tests for teacher evaluation in Chile (like the PAES or teacher certification exams), this type of question tests alignment between objectives and content. The content that most directly supports the objective about elected officials and representative democracy would be the method of appointing authorities. \newline I'll go with D. \\
  \textbf{Innospark-1T} & answer A \\
  & Let me analyze this question step by step. \newline The question asks: A teacher will start working with her 6th grade students based on the following plan: 6th Grade Unit: The political organization of Chile and the basic principles established by the Constitution. Learning Objective: Explain some fundamental elements of the democratic organization of Chile, including the division of state powers and their representation through elected officials (councilors, mayors, dep \newline Let me consider each option: \newline - Option A: Citizen participation \newline - Option B: Rights and duties \newline - Option C: Type of State \newline \textit{[\ldots\ truncated \ldots]} \newline All of these require understanding the method by which authorities are appointed (elected) in a representative democracy. \newline Based on this analysis, the answer is D. \\
  \specialrule{1.2pt}{8pt}{8pt}
  \multicolumn{2}{@{}l}{\textbf{chile\_460} \quad $\Delta = +100\%$} \\
  \textbf{Question} & To start the unit Design and Dissemination, in a 2nd Year secondary class, a teacher conducts a diagnostic assessment to determine the students' approach to design and its different branches. Once the assessment is applied, the teacher identifies that they have a confused notion about industrial design, as they exclusively relate it to the manufacturing of industrial parts on a large scale. What activity could the teacher propose to the students so they can understand the concept of industrial design? \newline A. Observe pieces of industrial design to analyze them, classify them, and make sketches of the exhibited pieces \newline B. Choose design pieces found in the classroom to investigate their origin, evolution, and relationship with their production context \newline C. Deconstruct an industrial design piece through photography, for the purpose of creating a new design piece that addresses current needs \newline D. Research iconic design pieces present in the history of industrial design, to subsequently create a design piece considering contextual needs \\
  \textbf{Gold} & B \\
  \midrule
  \textbf{Kimi-K2.6} & answer D \\
  & The user wants me to answer a multiple-choice question with ONLY the letter (A, B, C, or D). No explanation. \newline The scenario: 2nd Year secondary class, starting a unit on Design and Dissemination. Diagnostic assessment shows students think industrial design is exclusively about manufacturing industrial parts on a large scale. Teacher wants to propose an activity so they can understand the concept of industrial design. \newline I need to choose the best activity to correct this misconception. The misconception is that industrial design = manufacturing industrial parts at large scale. Industrial design is actually about designing products for mass production, considering aesthetics, functionality, ergonomics, etc. It's not just about mechanical/industrial parts, but everyday objects, furniture, appliances, etc. \newline Let's analyze each option: \newline A. Observe pieces of industrial design to analyze them, classify them, and make sketches of the exhibited pieces. \newline - This involves observation and classification. It might help them see what industrial design encompasses. But does it directly address the misconception about manufacturing parts? It shows them examples, which could broaden their view. However, it's somewhat passive. \newline \textit{[\ldots\ truncated \ldots]} \newline I think D is the intended answer because it explicitly mentions "history of industrial design" which is the standard way to teach the breadth of the concept, and "considering contextual needs" captures the human-centered aspect that distinguishes it from mere manufacturing. \newline Final decision: D. \\
  \textbf{Innospark-1T} & answer B \\
  & Let me analyze this question step by step. \newline The question asks: To start the unit Design and Dissemination, in a 2nd Year secondary class, a teacher conducts a diagnostic assessment to determine the students' approach to design and its different branches. Once the assessment is applied, the teacher identifies that they have a confused notion about industrial design, as they exclusivel \newline Let me consider each option: \newline - Option A: Observe pieces of industrial design to analyze them, classify them, and make sketches of the exhibited pieces \newline - Option B: Choose design pieces found in the classroom to investigate their origin, evolution, and relationship with their production context \newline - Option C: Deconstruct an industrial design piece through photography, for the purpose of creating a new design piece that addresses current needs \newline \textit{[\ldots\ truncated \ldots]} \newline The activity in Option B will help students discover that industrial design encompasses everyday objects they use, not just factory machinery or industrial components. \newline Based on this analysis, the answer is B. \\
  \specialrule{1.2pt}{8pt}{8pt}
  \bottomrule
\end{longtable}